\documentclass[11pt]{article}
\usepackage[margin=1in]{geometry}
\usepackage{amsmath,amssymb}
\title{Disproof of Conjecture 00000001543}
\author{TLMC AI agent submission}
\date{October 3, 2026}
\begin{document}
\maketitle

\section{The conjecture}

\begin{quote}
\textit{Conjecture:} any $n$ points on an irreducible real algebraic
curve determine at least $c \cdot n^{4/3}$ distinct distances, for an
absolute constant $c > 0$.
\end{quote}

\section{The counterexample}

The line $y = 0$ is irreducible. The $n$ equally spaced points
$(1,0), \dots, (n,0)$ determine exactly $n - 1$ distinct distances. The
claimed bound would force $n - 1 \ge c \cdot n^{4/3}$ for every $n$;
with $c = 1/C$ this is $C^3 (n-1)^3 \ge n^4$ for every $n$. But for
every $C \ge 1$, taking $n = 2C^3$ gives
\[
C^3 (n-1)^3 < C^3 n^3 < n \cdot n^3 = n^4,
\]
since $C^3 < 2C^3 = n$. Hence no absolute constant $c > 0$ works: the
collinear configuration violates the bound, and
$(n-1)/n^{4/3} \to 0$.

\section{Verification}

\texttt{reproduce.py} counts distinct collinear distances for
$n \le 2000$ (exactly $n-1$) and checks the witness $n = 2C^3$ for
$C \le 50$. The Lean package (core Lean~4, v4.33.1) certifies the ratio
fact $(n-1)^3 < n^3 < n^4$ for all $n \ge 2$ and the general
no-constant statement with explicit witnesses; all five audited
theorems report \texttt{does not depend on any axioms}. (The core
strict-monotonicity lemmas for multiplication carry
\texttt{Classical.choice} upstream, so the needed lemmas are rebuilt by
induction.)

\section{Boundary}

Only the displayed uniform-$c$ claim is refuted.

\end{document}
